\documentclass[10pt,a4paper]{amsart}
\usepackage[margin=19mm]{geometry}
\usepackage{amsmath,amssymb,mathtools}
\usepackage[scr=rsfs,cal=euler]{mathalfa}
\usepackage{xfrac}
\usepackage[unicode,hidelinks,bookmarks=false]{hyperref}
\newcommand{\ie}{i.e.\ }
\newcommand{\cf}{cf.\ }
\newcommand{\resp}{respectively\ }
\newcommand{\Subsection}{\subsection}
\providecommand{\nobreakdash}{\nobreak\hbox{-}}
\setlength{\emergencystretch}{2em}
\multlinegap0pt
\usepackage{enumerate}
\usepackage{amssymb}
\usepackage{mathtools}
\newtheorem{lemma}{Lemma}
\title{On Hypersurfaces of Space Forms that are Gradient Yamabe Solitons}
\author{Jahnabi Chakraborti, Anandateertha Mangasuli}
\date{Source: arXiv:2608.29298v1 (CC BY 4.0)}
\begin{document}
\maketitle

\noindent\emph{Source and licence.} On Hypersurfaces of Space Forms that are Gradient Yamabe Solitons. Version arXiv:2608.29298v1 (CC BY 4.0), distributed by arXiv under the Creative Commons Attribution 4.0 International licence, http://creativecommons.org/licenses/by/4.0/, as recorded in the arXiv licence metadata for this exact version and shown on the abstract page https://arxiv.org/abs/2608.29298v1. Full licence notice: \texttt{licenses/arxiv-2608.29298-CC-BY-4.0.txt}.\par\smallskip

\noindent\textit{Editorial scope.} Counted unit: the article's lemma lemma3 (source0.tex lines 737--739). The article's complete original proof of it is delivered with it (source0.tex lines 740--752, one \texttt{proof} environment of 13 lines). No mathematics is written, repaired or re-derived: the statement and the proof are the article's own lines, in the article's own notation, and no retained line is substituted or re-flowed.  Retained uncounted as the source-local context the proof needs: lines 693--695; lines 703--705.  Nothing was omitted from any retained span, so the package carries no omission record.  No retained line contains a citation, so the wrapper carries no bibliography and no bibliography entry.  No retained line contains a figure environment, an image-inclusion command or drawing code, so no image is delivered and the package declares no asset dependency.  Editorial formatting only, in addition: an AMS \texttt{amsart} preamble that restates the article's own declarations and macros used by the retained lines, namely the environments \texttt{lemma} on separate counters, together with the standard packages those lines need.  The retained environments print in this document as Lemma 1 in order of appearance, whereas the article prints the same items with the numbering of its own full document.  The complete native source of the article as distributed in the arXiv e-print for this exact version is bundled unchanged as \texttt{sources/208839/source0.tex}.
\begin{equation}\label{FE4}
    R^{\Sigma}=(n-2)(n-3)C+(n-2)^2 \left|\vec{H}^{\Sigma}_{\overline{M}}\right|^2-\left|B^{\Sigma}_{\overline{M}}\right|^2.
\end{equation}

\begin{equation}\label{FE6}
    L(u)=\sum_{i, j} \left\{\big(n-2)\overline{g}\left(\vec{H}^{\Sigma}_{\overline{M}}, \gamma \right)\delta_{ij}-B_{ij}\right\} u_{ij},
\end{equation}

\begin{lemma}\label{lemma3}
    If the scalar curvature of $\Sigma$, $R^{\Sigma}>(n-2)(n-3) C$ then the linear operator $L$ is elliptic.
\end{lemma}

\begin{proof}
    Since $R^{\Sigma}>(n-2)(n-3) C$ it follows from \eqref{FE4} that 
    \begin{equation}\label{FE21}
        \left|\vec{H}^{\Sigma}_{\overline{M}}\right|^2>0
    \end{equation}and
    \begin{equation} \label{FE5}
        (n-2)^2\left|\vec{H}^{\Sigma}_{\overline{M}}\right|^2\ge {k_i}^2,
    \end{equation}
    where $k_i$'s are the eigenvalues of the self-adjoint linear map associated with the bilinear form $B_{\gamma}(X, Y)= \overline{g}\left(B^{\Sigma}_{\overline{M}}(X, Y), \gamma\right)$.\\
    Now, $\left(\vec{H}^{\Sigma}_{\overline{M}}\right)^2=\left\{\overline{g}\left(\vec{H}^{\Sigma}_{\overline{M}}, \gamma \right)\right\}^2$. Therefore, \eqref{FE21} implies $\overline{g}\left(\vec{H}^{\Sigma}_{\overline{M}}, \gamma \right)>0$ or $\overline{g}\left(\vec{H}^{\Sigma}_{\overline{M}}, \gamma \right)<0$. Given $(M^n, g)$ is an oriented hypersurface of $\overline{M}^n(C)$, we can choose the orientation such that $\overline{g}\left(\vec{H}^{\Sigma}_{\overline{M}}, \gamma \right)>0$.\\
    Hence, \eqref{FE5} implies that $0<(n-2)\overline{g}\left(\vec{H}^{\Sigma}_{\overline{M}}, \gamma \right)-k_i<2(n-2)\left|\vec{H}^{\Sigma}_{\overline{M}}\right|$, where $(n-2)\overline{g}\left(\vec{H}^{\Sigma}_{\overline{M}}, \gamma \right)-k_i$'s are the eigenvalues of the operator $L$ defined in \eqref{FE6}.
    %Since all the eigenvalues of the operator $L$ are positive, it is elliptic.
\end{proof}

\end{document}
